% GENERATED FILE. Edit canonical lessons, then run npm run generate:books.
% canonical-source-hash: fnv1a64:182ff7350bb24df6
% canonical-lessons: TA-C369-camukam, TA-C369-potumakkal, TA-C369-makkaltokai, TA-C369-kutimakan, TA-C369-vakaippatuttu

\chapter{Society and Its People}
\label{ch:ta-society-and-its-people}
\hlchaptermodality{369}

\begin{chapteropening}
\textbf{By the end of this chapter:} \emph{I can talk about society, the public and citizens.}

Complete the last lesson of chapter 369: I can talk about society, the public and citizens.
\end{chapteropening}

\section[\texorpdfstring{\ta{சமூகம்} (camūkam)}{சமூகம் (camūkam)}]{\ta{சமூகம்} (camūkam) — society; a community}

\label{lesson:TA-C369-camukam}



\begin{quote}
Before the new one: say the Tamil for a seminar, then the Tamil for a question.
\end{quote}

\subsection*{You'll want to know: \ta{சமூகம்}}
\textbf{\ta{சமூகம்}} — \emph{camūkam} — \textquotedblleft{}society; a community\textquotedblright{}.

\begin{quote}
\ta{ஒவ்வொருவரும்} \ta{சமூகத்துக்கு} \ta{உதவ} \ta{வேண்டும்}.
\end{quote}

\begin{quote}
\emph{ovvoruvarum camūkattukku utava vēṇṭum.}
\end{quote}

\begin{quote}
Everyone should help society.
\end{quote}

\ta{சமூகத்துக்கு}, to society, with the oblique -\ta{அத்து}-. Before another noun, \ta{சமூக} means social.

\begin{cousinweb}
\textbf{\ta{சமூகம்}} (\emph{camūkam}) is borrowed from Sanskrit \emph{samūha}, \textquotedblleft{}heap, collection\textquotedblright{} (CDIAL 13241), whose Pali and Prakrit forms mean \textquotedblleft{}multitude\textquotedblright{}: people taken together.
\end{cousinweb}

\begin{grammarlens}[title={What you've built}]
One more word for giving reasons and opinions.
\end{grammarlens}

\subsection*{Your turn}
\begin{itemize}
\raggedright
  \item \emph{Say it:} \emph{camūkam}
  \item \emph{Say it:} \emph{camūkam}, once more
  \item \emph{Say it:} \emph{ovvoruvarum camūkattukku utava vēṇṭum.} — Everyone should help society.
  \item \emph{Recall:} say the Tamil for a seminar, then the Tamil for a question, then say \emph{camūkam} again
\end{itemize}

\begin{tcolorbox}[breakable,colback=teal!4,colframe=teal!35!black,title={Before you move on}]
What does \ta{சமூகம்} mean? (\textquotedblleft{}Society; a community\textquotedblright{}.) Say it once more.
\end{tcolorbox}
\section[\texorpdfstring{\ta{பொதுமக்கள்} (potumakkaḷ)}{பொதுமக்கள் (potumakkaḷ)}]{\ta{பொதுமக்கள்} (potumakkaḷ) — the public}

\label{lesson:TA-C369-potumakkal}



\begin{quote}
Before the new one: say the Tamil for a question, then the Tamil for society.
\end{quote}

\subsection*{You'll want to know: \ta{பொதுமக்கள்}}
\textbf{\ta{பொதுமக்கள்}} — \emph{potumakkaḷ} — \textquotedblleft{}the public\textquotedblright{}.

\begin{quote}
\ta{இந்தப்} \ta{பூங்கா} \ta{பொதுமக்களுக்குத்} \ta{திறந்திருக்கிறது}.
\end{quote}

\begin{quote}
\emph{intap pūṅkā potumakkaḷukkut tiṟantirukkiṟatu.}
\end{quote}

\begin{quote}
This park is open to the public.
\end{quote}

\ta{பொது}, common + \ta{மக்கள்}, people.

\begin{cousinweb}
\textbf{\ta{பொதுமக்கள்}} (\emph{potumakkaḷ}) is native: \emph{potu}, \textquotedblleft{}common, general, public\textquotedblright{} (DEDR 4507), and \emph{makkaḷ}, \textquotedblleft{}human beings, people\textquotedblright{} (DEDR 4616).
\end{cousinweb}

\begin{grammarlens}[title={What you've built}]
One more word for giving reasons and opinions.
\end{grammarlens}

\subsection*{Your turn}
\begin{itemize}
\raggedright
  \item \emph{Say it:} \emph{potumakkaḷ}
  \item \emph{Say it:} \emph{potumakkaḷ}, once more
  \item \emph{Say it:} \emph{intap pūṅkā potumakkaḷukkut tiṟantirukkiṟatu.} — This park is open to the public.
  \item \emph{Recall:} say the Tamil for a question, then the Tamil for society, then say \emph{potumakkaḷ} again
\end{itemize}

\begin{tcolorbox}[breakable,colback=teal!4,colframe=teal!35!black,title={Before you move on}]
What does \ta{பொதுமக்கள்} mean? (\textquotedblleft{}The public\textquotedblright{}.) Say it once more.
\end{tcolorbox}
\section[\texorpdfstring{\ta{மக்கள்தொகை} (makkaḷtokai)}{மக்கள்தொகை (makkaḷtokai)}]{\ta{மக்கள்தொகை} (makkaḷtokai) — population}

\label{lesson:TA-C369-makkaltokai}



\begin{quote}
Before the new one: say the Tamil for society, then the Tamil for the public.
\end{quote}

\subsection*{You'll want to know: \ta{மக்கள்தொகை}}
\textbf{\ta{மக்கள்தொகை}} — \emph{makkaḷtokai} — \textquotedblleft{}population\textquotedblright{}.

\begin{quote}
\ta{இந்த} \ta{நகரத்தின்} \ta{மக்கள்தொகை} \ta{வேகமாக} \ta{வளர்கிறது}.
\end{quote}

\begin{quote}
\emph{inta nakarattiṉ makkaḷtokai vēkamāka vaḷarkiṟatu.}
\end{quote}

\begin{quote}
The population of this city is growing fast.
\end{quote}

\ta{மக்கள்}, people, + a word for a total: the head count.

\begin{cousinweb}
\textbf{\ta{மக்கள்தொகை}} (\emph{makkaḷtokai}) is native: \emph{makkaḷ}, \textquotedblleft{}people\textquotedblright{} (DEDR 4616), and \emph{tokai}, \textquotedblleft{}sum, total\textquotedblright{} (DEDR 3476): the count of the people.
\end{cousinweb}

\begin{grammarlens}[title={What you've built}]
One more word for giving reasons and opinions.
\end{grammarlens}

\subsection*{Your turn}
\begin{itemize}
\raggedright
  \item \emph{Say it:} \emph{makkaḷtokai}
  \item \emph{Say it:} \emph{makkaḷtokai}, once more
  \item \emph{Say it:} \emph{inta nakarattiṉ makkaḷtokai vēkamāka vaḷarkiṟatu.} — The population of this city is growing fast.
  \item \emph{Recall:} say the Tamil for society, then the Tamil for the public, then say \emph{makkaḷtokai} again
\end{itemize}

\begin{tcolorbox}[breakable,colback=teal!4,colframe=teal!35!black,title={Before you move on}]
What does \ta{மக்கள்தொகை} mean? (\textquotedblleft{}Population\textquotedblright{}.) Say it once more.
\end{tcolorbox}
\section[\texorpdfstring{\ta{குடிமகன்} (kuṭimakaṉ)}{குடிமகன் (kuṭimakaṉ)}]{\ta{குடிமகன்} (kuṭimakaṉ) — a citizen}

\label{lesson:TA-C369-kutimakan}



\begin{quote}
Before the new one: say the Tamil for the public, then the Tamil for population.
\end{quote}

\subsection*{You'll want to know: \ta{குடிமகன்}}
\textbf{\ta{குடிமகன்}} — \emph{kuṭimakaṉ} — \textquotedblleft{}a citizen\textquotedblright{}.

\begin{quote}
\ta{ஒவ்வொரு} \ta{குடிமகனுக்கும்} \ta{சில} \ta{கடமைகள்} \ta{இருக்கின்றன}.
\end{quote}

\begin{quote}
\emph{ovvoru kuṭimakaṉukkum cila kaṭamaikaḷ irukkiṉṟaṉa.}
\end{quote}

\begin{quote}
Every citizen has some duties.
\end{quote}

Plural: \ta{குடிமக்கள்}, \emph{kuṭimakkaḷ}.

Plural \ta{குடிமக்கள்}; a woman citizen is \ta{குடிமகள்}.

\begin{cousinweb}
\textbf{\ta{குடிமகன்}} (\emph{kuṭimakaṉ}) is native: \emph{kuṭi}, \textquotedblleft{}house, family, town\textquotedblright{} (DEDR 1655), and \emph{makaṉ}, \textquotedblleft{}son, person\textquotedblright{} (DEDR 4616): a son of the town.
\end{cousinweb}

\begin{grammarlens}[title={What you've built}]
One more word for giving reasons and opinions.
\end{grammarlens}

\subsection*{Your turn}
\begin{itemize}
\raggedright
  \item \emph{Say it:} \emph{kuṭimakaṉ}
  \item \emph{Say it:} \emph{kuṭimakaṉ}, once more
  \item \emph{Say it:} \emph{ovvoru kuṭimakaṉukkum cila kaṭamaikaḷ irukkiṉṟaṉa.} — Every citizen has some duties.
  \item \emph{Recall:} say the Tamil for the public, then the Tamil for population, then say \emph{kuṭimakaṉ} again
\end{itemize}

\begin{tcolorbox}[breakable,colback=teal!4,colframe=teal!35!black,title={Before you move on}]
What does \ta{குடிமகன்} mean? (\textquotedblleft{}A citizen\textquotedblright{}.) Say it once more.
\end{tcolorbox}
\section[\texorpdfstring{\ta{வகைப்படுத்து} (vakaippaṭuttu)}{வகைப்படுத்து (vakaippaṭuttu)}]{\ta{வகைப்படுத்து} (vakaippaṭuttu) — to classify; to sort}

\label{lesson:TA-C369-vakaippatuttu}



\begin{quote}
Before the new one: say the Tamil for population, then the Tamil for a citizen.
\end{quote}

\subsection*{You'll want to know: \ta{வகைப்படுத்து}}
\textbf{\ta{வகைப்படுத்து}} — \emph{vakaippaṭuttu} — \textquotedblleft{}to classify; to sort\textquotedblright{}.

\begin{quote}
\ta{அவள்} \ta{புத்தகங்களைப்} \ta{பாடம்} \ta{வாரியாக} \ta{வகைப்படுத்தினாள்}.
\end{quote}

\begin{quote}
\emph{avaḷ puttakaṅkaḷaip pāṭam vāriyāka vakaippaṭuttiṉāḷ.}
\end{quote}

\begin{quote}
She sorted the books by subject.
\end{quote}

Past (I …): \ta{வகைப்படுத்தினேன்}, \emph{vakaippaṭuttiṉēṉ}.

What is sorted takes -\ta{ஐ}; \ta{வாரியாக}, by: \ta{பாடம்} \ta{வாரியாக}, subject by subject.

\begin{cousinweb}
\textbf{\ta{வகைப்படுத்து}} (\emph{vakaippaṭuttu}) is native: \emph{vakai}, \textquotedblleft{}division, branch, kind\textquotedblright{} (DEDR 5202), with \emph{paṭuttu}, the causative of \emph{paṭu}: to put into kinds.
\end{cousinweb}

\begin{grammarlens}[title={What you've built}]
One more word for giving reasons and opinions.
\end{grammarlens}

\subsection*{Your turn}
\begin{itemize}
\raggedright
  \item \emph{Say it:} \emph{vakaippaṭuttu}
  \item \emph{Say it:} \emph{vakaippaṭuttu}, once more
  \item \emph{Say it:} \emph{avaḷ puttakaṅkaḷaip pāṭam vāriyāka vakaippaṭuttiṉāḷ.} — She sorted the books by subject.
  \item \emph{Recall:} say the Tamil for population, then the Tamil for a citizen, then say \emph{vakaippaṭuttu} again
\end{itemize}

\begin{tcolorbox}[breakable,colback=teal!4,colframe=teal!35!black,title={Before you move on}]
What does \ta{வகைப்படுத்து} mean? (\textquotedblleft{}To classify; to sort\textquotedblright{}.) Say it once more.
\end{tcolorbox}
